\uuid{LT4V}
\titre{Quand une ReLU bloque l'apprentissage}
\niveau{L3}
\module{Analyse de données}
\chapitre{Réseaux de neurones}
\sousChapitre{Rétropropagation}
\theme{ReLU inactive, gradient, initialisation}
\auteur{Maxime Nguyen}
\datecreate{2026-09-26}
\organisation{AMSCC}
\difficulte{2}

\contenu{
\texte{On considère un réseau à une entrée, un neurone caché et une sortie linéaire :
\[
h(x)=\operatorname{ReLU}(wx+b),\qquad F(x)=v\,h(x),
\qquad \operatorname{ReLU}(t)=\max(0,t).
\]
Pour l'unique donnée d'entraînement $(x,t)=(1,2)$, la perte vaut
\[
E(w,b,v)=\frac12\bigl(F(1)-2\bigr)^2.
\]
On compare deux initialisations : $A=(-1,0,1)$ et $B=(1,0,1)$, dans l'ordre $(w,b,v)$.}

\begin{enumerate}
\item \question{Calculer $h(1)$, $F(1)$ et $E$ pour les initialisations $A$ et $B$.}
\indication{Vérifier le signe de $w+b$ dans chaque cas.}
\reponse{Pour $A$, $w+b=-1$, donc $h(1)=0$, $F(1)=0$ et $E=2$. Pour $B$, $w+b=1$, donc $h(1)=1$, $F(1)=1$ et $E=1/2$.}

\item \question{Lorsque $w+b\neq0$, calculer les trois dérivées partielles de $E$ en fonction de $F(1)$, $h(1)$ et de $\operatorname{ReLU}'(w+b)$. Les évaluer en $A$ et en $B$.}
\indication{Appliquer la règle de la chaîne ; la dérivée de ReLU vaut $0$ sur $]-\infty,0[$ et $1$ sur $]0,+\infty[$.}
\reponse{À l'entrée $x=1$, la règle de la chaîne donne
\[
\frac{\partial E}{\partial w}=(F(1)-2)v\,\operatorname{ReLU}'(w+b),\quad
\frac{\partial E}{\partial b}=(F(1)-2)v\,\operatorname{ReLU}'(w+b),\quad
\frac{\partial E}{\partial v}=(F(1)-2)h(1).
\]
Ainsi $\nabla E(A)=(0,0,0)$ et $\nabla E(B)=(-1,-1,-1)$, dans l'ordre $(w,b,v)$.}

\item \question{Effectuer une mise à jour simultanée par descente de gradient avec un pas $\eta=0{,}1$ depuis chaque initialisation. Pour le point obtenu depuis $B$, recalculer la sortie et la perte.}
\indication{Toutes les composantes d'une même mise à jour utilisent le gradient calculé aux paramètres initiaux.}
\reponse{Depuis $A$, le gradient nul laisse les paramètres inchangés. Depuis $B$, la mise à jour donne $(w,b,v)=(1{,}1,0{,}1,1{,}1)$. Alors $h(1)=1{,}2$, $F(1)=1{,}32$ et
\[
E=\frac12(1{,}32-2)^2=0{,}2312.
\]
La perte a diminué par rapport à $1/2$.}

\item \question{Expliquer pourquoi l'apprentissage est bloqué en $A$ pour cette donnée, malgré une perte non nulle.}
\indication{Relier la sortie nulle de la ReLU à ses dérivées et à la mise à jour des paramètres.}
\reponse{En $A$, la préactivation vaut $-1$ : la ReLU est inactive. Sa sortie et sa dérivée valent $0$, ce qui annule les trois composantes du gradient. La descente de gradient ne modifie donc aucun paramètre à partir de cette seule donnée, même si $E=2$. Une autre donnée ou une autre initialisation pourrait activer le neurone.}
\end{enumerate}
}
